\documentclass[11pt,a4paper]{article}
\usepackage[utf8]{inputenc}
\usepackage[T1]{fontenc}
\usepackage[french]{babel}
\usepackage{lmodern,amsmath,amssymb,amsfonts}
\usepackage{graphicx,booktabs,array}
\usepackage[margin=2.2cm,landscape]{geometry}
\usepackage{float,needspace,microtype,tikz}
\usepackage[colorlinks=true,linkcolor=blue!45!black,urlcolor=blue!45!black]{hyperref}
\usepackage[figure]{hypcap}
\usepackage{xcolor,xurl}
\setlength{\parskip}{4pt}
\setlength{\abovecaptionskip}{5pt}
\setlength{\belowcaptionskip}{0pt}
\hypersetup{pdftitle={ResoniCS -- Planning de la semaine},pdfauthor={Équipe ResoniCS}}
\pagestyle{empty}
\begin{document}
\section*{Organisation prévisionnelle du 2 au 6 novembre 2026}
Les tâches complètent les fichiers fournis dans support. Chaque modification porte sur un seul fichier, partagé entre les membres selon un passage de main. Les fiches individuelles détaillent les tâches et les contrôles de chaque jour.
\begin{figure}[H]
\centering
\setlength{\tabcolsep}{3pt}
\setlength{\extrarowheight}{2pt}
\newlength{\ganttTitleWidth}
\setlength{\ganttTitleWidth}{\dimexpr\textwidth-1.3cm-6cm-8cm-16\tabcolsep\relax}
\begin{tabular}{@{}p{1.3cm}p{\ganttTitleWidth}p{6cm}*{5}{>{\centering\arraybackslash}p{1.6cm}}@{}}
\toprule
\textbf{Repère} & \textbf{Tâche de l\textquotesingle énoncé} & \textbf{Responsable / contributions} & \shortstack{\textbf{Lundi}\\ 2 nov.} & \shortstack{\textbf{Mardi}\\ 3 nov.} & \shortstack{\textbf{\fontsize{9.5}{11}\selectfont Mercredi}\\ 4 nov.} & \shortstack{\textbf{Jeudi}\\ 5 nov.} & \shortstack{\textbf{Vendredi}\\ 6 nov.} \\
\midrule
% SOURCE_ID: J1.1
% SOURCE_TITLE: the state of the art ;
\textbf{J1.1} & the state of the art ; & \textbf{Olympia}\newline {\fontsize{9.5}{11}\selectfont Mathis} & \ganttbar{orange!80!black} &  &  &  &  \\
% SOURCE_ID: J1.2
% SOURCE_TITLE: the model to consider,
\textbf{J1.2} & the model to consider, & \textbf{Yann}\newline {\fontsize{9.5}{11}\selectfont Olympia, Joachim} & \ganttbar{blue!60!black} &  &  &  &  \\
% SOURCE_ID: J1.3
% SOURCE_TITLE: the plage of frequencies of the interest, dimensions, noise sources ;
\textbf{J1.3} & the plage of frequencies of the interest, dimensions, noise sources ; & \textbf{Mathis}\newline {\fontsize{9.5}{11}\selectfont Yann, Olympia} & \ganttbar{green!45!black} &  &  &  &  \\
% SOURCE_ID: J1.4
% SOURCE_TITLE: formulate the properties of the chosen model, what is known about its solution ;
\textbf{J1.4} & formulate the properties of the chosen model, what is known about its solution ; & \textbf{Joachim}\newline {\fontsize{9.5}{11}\selectfont ---} & \ganttbar{violet!70!black} &  &  &  &  \\
% SOURCE_ID: J1.5
% SOURCE_TITLE: formulate optimization problem ;
\textbf{J1.5} & formulate optimization problem ; & \textbf{Antoine}\newline {\fontsize{9.5}{11}\selectfont Joachim} & \ganttbar{red!65!black} &  &  &  &  \\
% SOURCE_ID: J1.6
% SOURCE_TITLE: Theoretical study : considering that one boundary is open-air (the wave go through without coming back) derive the energy over the form and compare the result with the analogous result with a reflexive boundary[^1].
\textbf{J1.6} & Theoretical study : considering that one boundary is open-air (the wave go through without coming back) derive the energy over the form and compare the result with the analogous result with a reflexive boundary\textsuperscript{1}. & \textbf{Joachim}\newline {\fontsize{9.5}{11}\selectfont Yann} & \ganttbar{violet!70!black} &  &  &  &  \\
% SOURCE_ID: J1.7
% SOURCE_TITLE: to choose porous material(s) and its quantity ;
\textbf{J1.7} & to choose porous material(s) and its quantity ; & \textbf{Mathis}\newline {\fontsize{9.5}{11}\selectfont Antoine} & \ganttbar{green!45!black} &  &  &  &  \\
% SOURCE_ID: J1.8
% SOURCE_TITLE: find numerically the corresponding complex coefficient(s) $\alpha(\omega)$ for the porous matherial(s) ; to compare the performances if several ;
\textbf{J1.8} & find numerically the corresponding complex coefficient(s) $\alpha(\omega)$ for the porous matherial(s) ; to compare the performances if several ; & \textbf{Mathis}\newline {\fontsize{9.5}{11}\selectfont Yann} & \ganttbar{green!45!black} &  &  &  &  \\
% SOURCE_ID: J1.9
% SOURCE_TITLE: find energy derivative over the distribution of the porous inclusions ;
\textbf{J1.9} & find energy derivative over the distribution of the porous inclusions ; & \textbf{Joachim}\newline {\fontsize{9.5}{11}\selectfont Yann} & \ganttbar{violet!70!black} &  &  &  &  \\
% SOURCE_ID: J1.10
% SOURCE_TITLE: formulate/consider the numerical gradient descent method.
\textbf{J1.10} & formulate/consider the numerical gradient descent method. & \textbf{Antoine}\newline {\fontsize{9.5}{11}\selectfont ---} & \ganttbar{red!65!black} &  &  &  &  \\
\midrule
% SOURCE_ID: J2.1
% SOURCE_TITLE: Study the energy on the chosen plage of frequencies for a fully absorbent wall of the plane shape and to define the local maxima ;
\textbf{J2.1} & Study the energy on the chosen plage of frequencies for a fully absorbent wall of the plane shape and to define the local maxima ; & \textbf{Olympia}\newline {\fontsize{9.5}{11}\selectfont Yann, Mathis} &  & \ganttbar{orange!80!black} &  &  &  \\
% SOURCE_ID: J2.2
% SOURCE_TITLE: Study the same question for different shapes of fully absorbent walls ;
\textbf{J2.2} & Study the same question for different shapes of fully absorbent walls ; & \textbf{Olympia}\newline {\fontsize{9.5}{11}\selectfont Yann, Mathis} &  & \ganttbar{orange!80!black} &  &  &  \\
% SOURCE_ID: J2.3
% SOURCE_TITLE: Compare the energy of the initial distribution of the porous material (this time the wall has not only absorbing parts but also reflective ones) with the energy for the fully absorbing wall.
\textbf{J2.3} & Compare the energy of the initial distribution of the porous material (this time the wall has not only absorbing parts but also reflective ones) with the energy for the fully absorbing wall. & \textbf{Olympia}\newline {\fontsize{9.5}{11}\selectfont Yann, Antoine} &  & \ganttbar{orange!80!black} &  &  &  \\
% SOURCE_ID: J2.4
% SOURCE_TITLE: Optimize numerically for one or several of the most problematic frequencies (at which there is a maximum of the energy) ;
\textbf{J2.4} & Optimize numerically for one or several of the most problematic frequencies (at which there is a maximum of the energy) ; & \textbf{Antoine}\newline {\fontsize{9.5}{11}\selectfont Yann, Olympia, Joachim} &  & \ganttbar{red!65!black} &  &  &  \\
\midrule
% SOURCE_ID: J3.1
% SOURCE_TITLE: Analyze the obtained results to compare to the efficient performances on the fixed range of frequencies ; to compare to a fully absorbing wall ;
\textbf{J3.1} & Analyze the obtained results to compare to the efficient performances on the fixed range of frequencies ; to compare to a fully absorbing wall ; & \textbf{Olympia}\newline {\fontsize{9.5}{11}\selectfont Yann} &  &  & \ganttbar{orange!80!black} &  &  \\
% SOURCE_ID: J3.2
% SOURCE_TITLE: to optimize the energy at once at several different frequencies ;
\textbf{J3.2} & to optimize the energy at once at several different frequencies ; & \textbf{Antoine}\newline {\fontsize{9.5}{11}\selectfont Yann, Olympia, Joachim} &  &  & \ganttbar{red!65!black} &  &  \\
% SOURCE_ID: J3.3
% SOURCE_TITLE: to study which quantity of the porous material gives the most optimal result ;
\textbf{J3.3} & to study which quantity of the porous material gives the most optimal result ; & \textbf{Mathis}\newline {\fontsize{9.5}{11}\selectfont Yann, Olympia, Antoine} &  &  & \ganttbar{green!45!black} &  &  \\
\midrule
% SOURCE_ID: J4.1
% SOURCE_TITLE: To continue the analysis of the performances ;
\textbf{J4.1} & To continue the analysis of the performances ; & \textbf{Olympia}\newline {\fontsize{9.5}{11}\selectfont Yann, Mathis, Joachim, Antoine} &  &  &  & \ganttbar{orange!80!black} &  \\
% SOURCE_ID: J4.2
% SOURCE_TITLE: To conclude which product to produce/develop.
\textbf{J4.2} & To conclude which product to produce/develop. & \textbf{Olympia}\newline {\fontsize{9.5}{11}\selectfont Yann, Mathis, Joachim, Antoine} &  &  &  & \ganttbar{orange!80!black} &  \\
\midrule
% SOURCE_ID: J5.1
% SOURCE_TITLE: To finalize everything
\textbf{J5.1} & To finalize everything & \textbf{Yann}\newline {\fontsize{9.5}{11}\selectfont Mathis, Olympia, Joachim, Antoine} &  &  &  &  & \ganttbar{blue!60!black} \\
% SOURCE_ID: J5.2
% SOURCE_TITLE: Preparation of the defense
\textbf{J5.2} & Preparation of the defense & \textbf{Yann}\newline {\fontsize{9.5}{11}\selectfont Mathis, Olympia, Joachim, Antoine} &  &  &  &  & \ganttbar{blue!60!black} \\
% SOURCE_ID: J5.3
% SOURCE_TITLE: Finalization of the rapport by the team.
\textbf{J5.3} & Finalization of the rapport by the team. & \textbf{Yann}\newline {\fontsize{9.5}{11}\selectfont Mathis, Olympia, Joachim, Antoine} &  &  &  &  & \ganttbar{blue!60!black} \\
\bottomrule
\end{tabular}

\caption{Répartition des travaux sur les cinq journées}
\end{figure}
\end{document}
